\begin{resumo}[Abstract]
This work evaluates the capability of five Large Language Models (LLMs) --- ChatGPT, Gemini, Claude, Meta AI, and DeepSeek --- to convert sequential C programs into parallel implementations using OpenMP. 
The evaluation was conducted on six programs from the PAMPAR benchmark, ranging from reductions with independent iterations to algorithms with complex data dependencies. 
Each model received the same standardized prompt and the sequential source code, in a single attempt, without iterative refinement. 
The generated programs were assessed for compilation, output correctness, and the speedup achieved on a shared-memory machine, using PAMPAR's reference implementations as a performance target. 
All generated programs compiled and produced output identical to the sequential version. 
In terms of performance, the frontier models approached the expert-written reference implementation, while the remaining models showed efficiency losses concentrated in the programs with more complex parallelization patterns. 
The results indicate that current LLMs are able to generate functionally correct parallel code for a wide variety of programs in a single attempt, but differ significantly in efficiency as parallelization complexity increases.

\vspace{\onelineskip}
\noindent
\textbf{Keywords}: Large Language Models. Automatic Parallelization. OpenMP. High-Performance Computing.
\end{resumo}